% ===========================================================================
\section{Methods}
\label{sec:methods}
% ===========================================================================

The framework consists of four main stages: calculation of the national agricultural
inventory; allocation of national totals to provinces; allocation to a \SI{0.05}{\degree}
grid and to administrative districts; and uncertainty, key-category, completeness and
quality-control analysis.

The inventory is predominantly Tier~1. Each category uses the equations in the 2006 IPCC
Guidelines, together with relevant updates from the 2019 Refinement, and applies them to
country-specific activity data. Selected livestock parameters provide partial Tier~2
characterisation where nationally appropriate values are available. These include typical
animal masses, the dairy-cattle share and manure-management-system shares.

Tier status is reported by category in Table~\ref{tab:tierstatus}. This is more informative
than assigning one tier to the entire inventory, because the level of country-specific
detail differs among categories.

All activity values, emission factors and parameters are stored in version-controlled data
files; values are not embedded directly in the calculation code. Greenhouse gases are
converted to carbon-dioxide equivalents using AR5 100-year global-warming potentials ---
28 for \CHfour{} and 265 for \NtwoO{} --- to match the Common Reporting Table convention.
AR6-equivalent totals are reported separately.

\begin{table}[!htbp]
\centering
\footnotesize
\setlength{\tabcolsep}{4pt}
\begin{tabular}{@{}lp{0.26\linewidth}p{0.34\linewidth}@{}}
\toprule
\textbf{Category} & \textbf{Method / data} & \textbf{Country-specific elements} \\
\midrule
Enteric fermentation & T1 method, CS data & species census, dairy split \\
Manure management & T1/2, CS data & animal mass, Nex, MMS, temp.\ zone \\
Managed soils & T1 method, CS data & fertiliser data, dry-climate EFs \\
Rice cultivation & T1 method, CS data & area, cultivation period \\
Field burning & T1 method, CS data & crop production, residue params \\
Urea application & T1 method, CS data & urea offtake; N series before 2024 \\
\bottomrule
\end{tabular}
\caption{Method-tier status by category. All categories use the Tier-1 (T1)
method equations on country-specific (CS) activity data; manure management is
partly Tier~2 through country-specific animal mass and management-system shares.
No emission factor is Tier~3 (Pakistani field measurement).}
\label{tab:tierstatus}
\end{table}

\subsection{National inventory}

\paragraph{Enteric fermentation and manure management}
Enteric methane is calculated using the Tier~1 form of IPCC Eq.~10.19,
\begin{equation}
\mathrm{CH_4^{ent}} \;=\; \sum_t N_t\, EF_t \times 10^{-6},
\end{equation}
where $N_t$ is the population of livestock category $t$ and $EF_t$ its annual per-head
methane emission factor.

Methane from manure management is calculated using IPCC Eq.~10.22 with temperature-dependent
factors, and direct nitrous oxide from manure management follows Eq.~10.25. Nitrogen
excretion is calculated as $N_{\mathrm{rate},t} \times TAM_t \times 365/1000$, where
$N_{\mathrm{rate},t}$ is the daily nitrogen-excretion rate and $TAM_t$ is typical animal
mass.

The calculation then distributes excreted nitrogen among manure-management systems. Dung
burned as household fuel is assigned to the energy sector and excluded from the agricultural
inventory to prevent double counting. Nitrogen deposited directly on pasture is assigned to
managed soils rather than to manure management.

\paragraph{Managed soils}
Direct soil nitrous oxide follows IPCC Eq.~11.1,
\begin{equation}
\mathrm{N_2O^{direct}} \;=\;
\left[(F_{\mathrm{SN}}+F_{\mathrm{ON}}+F_{\mathrm{CR}})\,EF_1
+ F_{\mathrm{PRP}}\,EF_{3,\mathrm{PRP}}\right]\frac{44}{28},
\end{equation}
where $F_{\mathrm{SN}}$, $F_{\mathrm{ON}}$, $F_{\mathrm{CR}}$ and $F_{\mathrm{PRP}}$ are
nitrogen inputs from synthetic fertiliser, organic amendments, crop residues and pasture
deposition respectively.

Nitrogen applied to rice is assigned the flooded-rice emission factor
$EF_{1,\mathrm{FR}}$. Indirect \NtwoO{} from volatilisation and leaching follows IPCC
Eqs.~11.9 and~11.10.

For direct \NtwoO{} from non-rice cropland, the reference retains the 2019 dry-climate value $EF_1=0.005$ as a conditional factor transfer. Aridity alone does not establish its applicability to irrigated cropland. Rice and grazing use their separate source formulations, and aggregate/climate alternatives remain structural sensitivities. Agreement with an external inventory is not used to select the factor.

\paragraph{Rice cultivation, residue burning and urea application}
Rice methane is calculated using IPCC Eqs.~5.1 and~5.2,
\begin{equation}
\mathrm{CH_4^{rice}} \;=\; \sum_i EF_c\, SF_{w,i}\, SF_{p,i}\, SF_{o,i}\, t_i\, A_i
\times 10^{-6},
\end{equation}
where $EF_c$ is the baseline daily emission factor, $SF_w$ represents the water regime,
$SF_p$ the pre-season water regime, $SF_o$ the effect of organic amendments, $t_i$ the
cultivation period and $A_i$ the harvested area.

Emissions from field burning follow the generic biomass-burning equation in the IPCC
Guidelines. Urea-application \COtwo{} follows Eq.~11.13. The urea mass has two bases, and
each year states which it used in \texttt{outputs/agriculture\_urea\_basis.csv}. For the
reference year it is the national urea offtake NFDC prints for the fiscal year itself. For
every earlier year, which has a total-nitrogen figure and no published product split, it is
reconstructed from that series by a constant chemical relation. The reconstruction was
formerly applied to the reference year as well, where it fell short of the printed offtake by
\SI{0.207}{\percent}; a measurement is preferred to an approximation of it wherever one
exists, and the approximation is retained, and labelled, where none does.

Liming emissions under category 3.G are not estimated: no national lime-use activity figure
has been sought. Liming is recorded as ``not estimated'' in the completeness
assessment.

\subsection{Provincial allocation}

National emission streams are allocated to the four provinces and a small residual region
using activity-specific administrative shares.

Enteric fermentation, manure-management emissions and pasture-deposited nitrogen use
provincial livestock-census shares. Rice methane and rice-specific fertiliser nitrogen use
provincial rice area. Residue-related emissions use cropland shares. Synthetic-fertiliser
\NtwoO{} and urea \COtwo{} use provincial fertiliser-offtake shares.

Applied manure requires a different treatment, because it is generated near livestock but
applied to cropland. Its provincial distribution is therefore derived by intersecting
cropland with livestock density at the grid-cell level and then aggregating the resulting
surface to provinces.

All provincial shares are normalised to sum to one. Provincial estimates therefore reproduce
the national control total exactly.

\subsection{Gridded allocation and mass-balance reconciliation}

Each provincial emission stream is allocated to a \SI{0.05}{\degree} grid using dasymetric
weighting. Let $E_{p,s}$ be the control total for emission stream $s$ in province $p$, let
$\Omega_p$ be the set of grid cells within that province, and let $w_{i,s}$ be the
spatial-proxy value for cell $i$. Emissions in each cell are calculated as
\begin{equation}
e_{i,s} \;=\; E_{p,s}\,\frac{w_{i,s}}{\sum_{j\in\Omega_p} w_{j,s}},
\qquad i\in\Omega_p .
\label{eq:reconcile}
\end{equation}
This construction guarantees that $\sum_{i\in\Omega_p} e_{i,s}=E_{p,s}$. The grid therefore
reconciles to the provincial total, and provincial totals reconcile to the national
inventory.

Spatial proxies are assigned by emission stream rather than by broad inventory category.
This distinction matters because a single category can contain several nitrogen or activity
streams with different geographies.

\GLW{} livestock density is used for enteric methane, manure-management emissions and
pasture-deposited nitrogen. Species densities are weighted by their per-head emission
factors, so that the resulting proxy represents the spatial distribution of emissions rather
than unweighted animal numbers.

MapSPAM paddy area is the declared proxy for rice methane and the rice-specific
fertiliser stream; the raster (2020 V2r2, hash-pinned) is held and consumed since
the spatial facility program, its cropland-fallback class deleted fail-closed,
and the own-year GCD-Rice 2023 classification ships beside it as a published
crosscheck and a priced district-level bracket (Section~\ref{sec:data}). FIRMS
active-fire detections are used for crop-residue burning. WorldCover cropland fraction is
used for synthetic fertiliser, crop-residue inputs to soils and urea application.

Managed soils are constructed by allocating each nitrogen input separately and then summing
the resulting stream-level surfaces. This prevents the geography of synthetic fertiliser
from being imposed on manure-derived or pasture-derived nitrogen.

Applied manure is allocated using cropland weighted by livestock density. This proxy
represents two physical constraints: manure is applied on agricultural land, but transport
costs generally keep it close to the herd that produces it.

Where a proxy has no positive values within a province, the corresponding emission stream is
distributed uniformly across that province. This fallback preserves the provincial control
total.

\subsection{District allocation and historical district panel}
\label{sec:methods-district}

The same proportional-allocation rule is used to distribute provincial totals among
administrative districts. The main output contains 116 districts across the four provinces.
Gilgit-Baltistan, Azad Jammu and Kashmir, and Islamabad Capital Territory are represented
together as a residual unit, producing 115 reporting units in total.

For most emission streams, district shares are obtained by summing the relevant spatial
proxy within each district and normalising those sums within the province.

Livestock categories receive a more direct treatment: livestock emissions are distributed
using district head-counts from the 2024 PBS Agricultural Census, every census district
resolved through the inventory's canonical district crosswalk (several census districts
sum onto one boundary polygon where the boundary file predates a split). The two census
units the boundary file does not draw, Cholistan Area and Larkana District, keep their own
share as named rows of the district panel and fall to the province surface in the grid; a
polygon the census does not reach is a build failure, not a proxy.

The gridded and district products are parallel allocations of the same provincial controls.
They are not nested representations of one another. The \SI{0.05}{\degree} grid uses
continuous gridded proxies, whereas the district table substitutes census head-counts for
livestock where possible. Both products reconcile to the same provincial and national
totals, but summing grid cells within a district will not necessarily reproduce the
census-controlled district estimate. Formally,
$\sum_{i\in\Omega_p}e_{i,s}=E_{p,s}$ for the grid and $\sum_{d\in p}E_{d,s}=E_{p,s}$ for the
district product, and both satisfy $\sum_p E_{p,s}=E_s$.

No consistent multi-year district activity dataset is available. The historical district
panel therefore holds the audited 2024 district shares fixed and rescales them to each year's
national category total for 1990--2024.

The district values consequently follow the national time series. They do not represent
observed changes in the geographic distribution of activity over time. This assumption is
evaluated in Section~\ref{sec:limitations}.

The exploratory land-sector extension in Supplementary Note~S1 is excluded from the
agricultural totals, benchmarking, key-category assessment and uncertainty analysis.

\subsection{Uncertainty analysis}

Uncertainty is estimated using two propagation approaches applied to the same set of
activity-data and emission-factor uncertainties.

IPCC Approach~1 propagates uncertainty analytically. Uncertainties in multiplicative terms
are combined in quadrature, while uncertainties in additive components are combined using
their contributions to the total.

The Monte-Carlo analysis samples the \emph{components}: each engine row is drawn from a
mean-preserving lognormal at its own declared width---appropriate for non-negative
emissions with potentially right-skewed uncertainty---with one shared deviate per
declared correlation group, so rows split apart for allocation reasons cannot average
their common error away. The simulation uses $10^4$ iterations, and the category and
sector totals are the \emph{sums of the sampled rows} in every realisation rather than
sampled quantities in their own right.

The two approaches test the sensitivity of the reported interval to the form of propagation.
They do not provide independent estimates, because both use the same underlying uncertainty
assumptions. Agreement between them therefore indicates consistency between analytical
quadrature and lognormal aggregation.

Managed-soils components that share an emission factor are treated as perfectly correlated.
This applies to separate volatilisation and leaching streams that appear as distinct rows
because they use different spatial proxies. Treating them as independent would artificially
reduce uncertainty when one physical stream is divided into several reporting components.

Other within-category and cross-category correlations are not modelled \emph{by this
sector-scoped stream}, and the qualifier matters because a second engine treats two of them
differently. This sector's Monte Carlo samples component totals, which cannot express
covariance between the categories its stages couple; the NATIONAL run instead samples a
parameter graph, and there two physical couplings are carried explicitly---fertiliser N as
one audited series driving both the soils components and the urea term, and each species'
head-count shared between its enteric and manure components. Their effect at national scale
is about a hundredth of a percentage point, in the widening direction. Supplementary
Note~S3 gives that treatment; the limitation of the sector-scoped stream is the one stated
here. The reported
national interval therefore reflects the documented parameter uncertainties under this
correlation structure. Factor-level uncertainties are reported in
Table~\ref{tab:ef-full}, category distributions in \ref{app:unc}, and full input
provenance in Supplementary Note~S3.

\subsection{Key categories, completeness and quality control}

Key categories are identified using the IPCC Approach~1 level and trend assessments.

The level assessment ranks categories by their contribution to the absolute national total
in the reference year. The trend assessment uses
$T_x = L_{x,0}\,\lvert g_x - g_{\mathrm{tot}}\rvert$, where $L_{x,0}$ is category $x$'s
share in the base year, $g_x$ is its growth since 1990 and $g_{\mathrm{tot}}$ is the growth
of the total agricultural inventory. The metric therefore identifies categories that are
both important in the base year and whose trajectory differs from the sector average. It
should not be interpreted simply as a category growth rate.

Categories are accumulated until the \SI{95}{\percent} threshold is reached. A category is
classified as key if it qualifies under either the level or the trend assessment.

Inventory coverage is documented through a completeness table. Every \AFOLU{}
category-and-gas combination receives an IPCC notation key: estimated, not estimated, not
occurring, or included elsewhere. This makes known gaps, exclusions and cross-sector
transfers explicit.

Quality control is implemented through automated tests covering units, source traceability,
mass-balance reconciliation, benchmark ranges, and consistency among national, provincial,
grid, district and historical outputs. Quality assurance is supported through comparison
with FAOSTAT, EDGAR, Climate TRACE and independent district statistics. Changes to activity
data, emission factors or methods trigger recalculation of the complete time series.
